% ===== PRÉAMBULE CANONIQUE — à coller VERBATIM en tête de CHAQUE .tex =====
\documentclass[11pt,a4paper]{article}
\usepackage[utf8]{inputenc}\usepackage[T1]{fontenc}\usepackage{lmodern}
\usepackage[french]{babel}
\usepackage{amsmath,amssymb,mathtools}\usepackage{icomma}
\usepackage{siunitx}\sisetup{locale=FR}  % \qty{}{}, \si{}, \ang{}
\usepackage{xcolor}\usepackage{colortbl}
\usepackage{tikz}\usepackage{pgfplots}\pgfplotsset{compat=1.18}
\usepackage[siunitx,europeanresistors]{circuitikz} % circuits (élec.)
\usepackage[most]{tcolorbox}\usepackage{enumitem}\usepackage{array,booktabs}
\usepackage[a4paper,margin=2.2cm]{geometry}
\usetikzlibrary{arrows.meta,calc,angles,quotes,positioning,patterns,%
  backgrounds,shapes.geometric,decorations.pathmorphing,decorations.markings,intersections}
\definecolor{primary}{RGB}{222,49,99}\definecolor{primarydark}{RGB}{150,22,55}
\definecolor{defcol}{RGB}{0,114,178}\definecolor{methcol}{RGB}{52,92,148}
\definecolor{excol}{RGB}{0,135,90}\definecolor{attncol}{RGB}{200,40,55}
\tcbset{boxbase/.style={enhanced,breakable,boxrule=0.9pt,arc=2.5pt,
  left=3mm,right=3mm,top=2.3mm,bottom=2.3mm,fonttitle=\bfseries,coltitle=white}}
% --- boîtes TUTORIEL (noms EXACTS, ne pas renommer) ---
\newtcolorbox{cle}[1][]{boxbase,colback=defcol!6,colframe=defcol,
  title={\textsc{À retenir}\ifx\relax#1\relax\relax\else\textmd{ : \itshape #1}\fi}}
\newtcolorbox{methode}[1][]{boxbase,colback=methcol!7,colframe=methcol,sharp corners,
  title={\textsc{Méthode}\ifx\relax#1\relax\relax\else\textmd{ --- \itshape #1}\fi}}
\newtcolorbox{exemplebox}[1][]{boxbase,colback=excol!5,colframe=excol,
  title={\textsc{Exemple}\ifx\relax#1\relax\relax\else\textmd{ : \itshape #1}\fi}}
\newtcolorbox{piege}[1][]{boxbase,colback=attncol!6,colframe=attncol,
  title={\textsc{Piège}\ifx\relax#1\relax\relax\else\textmd{ : \itshape #1}\fi}}
\newtcolorbox{remarque}[1][]{boxbase,colback=black!4,colframe=black!45,
  title={\textsc{Remarque}\ifx\relax#1\relax\relax\else\textmd{ : \itshape #1}\fi}}
% --- boîtes EXERCICE / CORRIGÉ (noms EXACTS, auto-comptées) ---
\newtcolorbox[auto counter]{exo}[1][]{boxbase,colback=primary!4,colframe=primary,
  title={\textsc{Exercice}~\thetcbcounter\ifx\relax#1\relax\relax\else\textmd{ --- \itshape #1}\fi}}
\newtcolorbox[auto counter]{corrige}[1][]{boxbase,colback=excol!4,colframe=excol,
  title={\textsc{Corrigé}~\thetcbcounter\ifx\relax#1\relax\relax\else\textmd{ --- \itshape #1}\fi}}
\newcommand{\vv}[1]{\vec{#1}}\newcommand{\ds}{\displaystyle}
\newcommand{\R}{\mathbb{R}}\newcommand{\N}{\mathbb{N}}\newcommand{\Z}{\mathbb{Z}}
% (un \usepackage{fancyhdr}+en-tête est possible ; il est ignoré côté web)
% =========================================================================
\begin{document}
\sisetup{per-mode=symbol}

\begin{exo}[vrai ou faux sur un graphe de vitesse]
Le graphe donne la vitesse d'un mobile $M$ le long de l'axe $Ox$.
\begin{center}
\begin{tikzpicture}
\begin{axis}[width=10cm,height=5cm,axis lines=middle,xlabel={$t\ (\si{\second})$},ylabel={$v\ (\si{\metre\per\second})$},
  xmin=0,xmax=7.6,ymin=-4,ymax=7.5,xtick={0,1,2,3,4,5,6,7},ytick={-3,0,3,6},grid=both,
  grid style={black!12},clip=false]
  \addplot[primary,line width=1.3pt] coordinates {(0,0)(2,6)(4,6)(7,-3)};
  \addplot[black,only marks,mark=*,mark size=1.2pt] coordinates {(6,0)};
\end{axis}
\end{tikzpicture}
\end{center}
Indique si chaque affirmation est vraie ($V$) ou fausse ($F$), puis compte les affirmations vraies.
\begin{enumerate}[label=\alph*)]
  \item Entre $t=0$ et $t=\qty{2}{\second}$, le mouvement est uniformément accéléré.
  \item Entre $t=2$ et $t=\qty{4}{\second}$, le mobile est à l'arrêt.
  \item À $t=\qty{6}{\second}$, le mobile change de sens.
  \item Entre $t=4$ et $t=\qty{7}{\second}$, l'accélération vaut \qty{-3}{\metre\per\second\squared}.
  \item La distance parcourue entre $t=0$ et $t=\qty{2}{\second}$ vaut \qty{6}{\metre}.
\end{enumerate}
\end{exo}

\begin{exo}[décrypter une équation horaire parabolique]
Un mobile a pour position $x(t)=-t^2+6t-5$ (en \si{\metre}, $t\ge0$ en \si{\second}).
\begin{enumerate}[label=\alph*)]
  \item Donne $v(t)$ et $a(t)$. De quel type de mouvement s'agit-il ?
  \item À quel instant la vitesse s'annule-t-elle ? Que représente ce point sur la trajectoire ?
  \item À quels instants le mobile passe-t-il par l'origine $x=0$ ?
  \item Dans quel sens se déplace-t-il avant, puis après l'instant trouvé au b) ?
\end{enumerate}
\end{exo}

\begin{exo}[déplacement contre distance parcourue]
La vitesse d'un mobile sur $Ox$ est $v(t)=-2t+4$ (\si{\metre\per\second}), $t$ en \si{\second}.
\begin{enumerate}[label=\alph*)]
  \item Écris l'intégrale donnant le \emph{déplacement} entre $t=0$ et $t=\qty{3}{\second}$, puis calcule-le.
  \item À quel instant le mobile change-t-il de sens ?
  \item En déduis la \emph{distance réellement parcourue} entre $t=0$ et $t=\qty{3}{\second}$.
\end{enumerate}
\end{exo}

\begin{exo}[affirmations de cours à trancher]
Pour chaque proposition, réponds par $V$ ou $F$ et justifie en une ligne. Combien sont exactes ?
\begin{enumerate}[label=\alph*)]
  \item Une accélération constante et non nulle impose une vitesse qui varie linéairement avec le temps.
  \item Si la vitesse instantanée est nulle à un instant, alors l'accélération y est forcément nulle.
  \item Le déplacement entre deux instants est l'aire algébrique sous la courbe $v(t)$.
  \item Un mobile dont l'accélération est négative ralentit toujours.
  \item La vitesse moyenne d'un aller-retour revenant au point de départ est nulle.
\end{enumerate}
\end{exo}

\end{document}
